\subsection{{Senior AI/ML Engineer (Personalization) \hfill May 2025 - Present}}  
\subtext{Target Corporation \hfill Minneapolis, MN}
\begin{zitemize}
\item Fast-tracked the productionization of a bundling algorithm in Vertex AI leading to projected lift of ~\textbf{\$20M attributed annual revenue} and \textbf{8\% increase in average cart value} getting deployed in 1 quarter after demo.    
\item Implemented the first internal tool to aid in \textbf{Blue-Green deployment} on AlloyDB, fully eliminating table downtime for updates \textbf{from 2 hours to sub-millisecond}. 
\item Created and distributed 5 \textbf{agentic skills and workflows} across 4 teams (total 50 engineers and data scientists) to boost developer efficiency with stacked PR creation, code review, and migration of entire repos. Supplemented these agentic skills with a presentation on context engineering so engineers know how to use tokens efficiently to achieve an output.
\item \textbf{Led} organization-wide response to a GCP service outage, using agentic orchestration to deploy code changes across \textbf{100+ repositories} and coordinating \textbf{4 teams} to contain impact before the weekend; established the standard for large-scale fleet deployments at Target. Continued to build out service for fleet deployment from this incident.
\item Built an \textbf{AI-powered triage assistant} integrated with \textbf{Slack}, \textbf{GitHub}, and \textbf{Vertex AI}, allowing teams to view error logs and restart the pipeline from their phone, leading to faster on-call triaging. 
\item Led migration of \textbf{10+ repositories} across Items and Content personalization teams from on-prem to \textbf{Vertex AI}.
\item Led the refactor and deprecation of an outdated and unsupported graph library and fully replaced it with \textbf{PyTorch Geometric}, while ensuring the model behavior is \textbf{1:1 identical} by utilizing \textbf{unit tests, integration tests, and holdout validation}, ensuring no statistically significant difference between the models. 
\end{zitemize}

\subsection{{AI/ML Engineer \hfill July 2024 - May 2025}}  
\subtext{NCR Voyix \hfill Atlanta, GA}  

\begin{zitemize}
\item Built scalable ML models for time-series prediction, improving the accuracy of KPI predictions from 16\% to 10\% error over a 3-week forecast using \textbf{Kubeflow}, \textbf{Google Cloud}, and \textbf{BigQuery} to allow customers to make data-driven business decisions.  
\item Built an end-to-end ML pipeline automating \textbf{data ingestion}, \textbf{preprocessing}, \textbf{feature engineering}, training, and inference for millions of transactions serving customer in a batch process while keeping costs to a minimum.
\item Automated model retraining, hyperparameter tuning, and deployment with \textbf{Kubeflow Pipelines}, \textbf{Katib}, and \textbf{GitHub Actions}, eliminating a tedious 2-hour weekly manual process.  
\item Integrated \textbf{CI/CD} for ML models, adding automated testing and validation, reducing failures, and allowing per-PR deployments.  

\end{zitemize}  



%====================
% EXPERIENCE A
%====================
\subsection{{Research Assistant \hfill Oct, 2020 --- June 2024}}
\begin{zitemize}
\item \textbf{FOCUS: Fast Observation and Collection of Soil \href{https://scholarworks.gsu.edu/cs_theses/115/}{(Primary Author, Master’s Thesis at GSU)}}
\item \textbf{Camera-Assisted Visible Light Communication \href{https://ieeexplore.ieee.org/abstract/document/9973607}{(Primary Author, IEEE REUNS 2022)}}
\end{zitemize}


%====================
% EXPERIENCE B
%====================

%====================
% EXPERIENCE E
%====================
%\subsection{{ROLE / PROJECT E \hfill MMM YYYY --- MMM YYYY}}
%\subtext{company E \hfill somewhere, state}
%\begin{zitemize}
%\item In lobortis libero consectetur eros vehicula, vel pellentesque quam fringilla.
%\item Ut malesuada purus at mi placerat dapibus.
%\item Suspendisse finibus massa eu nisi dictum, a imperdiet tellus convallis.
%\item Nam feugiat erat vestibulum lacus feugiat, efficitur gravida nunc imperdiet.
%\item Morbi porta lacus vitae augue luctus, a rhoncus est sagittis.
%\end{zitemize}
